%------------------------------------------------------------------------------%
\begin{exercises}

%------------------------------------------------------------------------------%
\begin{exercise}
Use substituição trigonométrica para calcular as integrais
\begin{mathlist}
  \item \int_0^1 x^3 \sqrt{1-x^2} dx
  \item \int_{\sqrt{2}}^2 \frac{1}{t^3(\sqrt{t^2-1})} dt
  \item \int \frac{dx}{\sqrt{x^2+16}} dx
  \item \int_0^{a} x^2 \sqrt{a^2-x^2}dx
  \item \int \frac{t^5}{\sqrt{t^2+2}} dt
  \item \int_{\frac{\sqrt{2}}{3}}^{\frac{2}{3}} \frac{dx}{x^5\sqrt{9x^2-1}} dx
  \item \int  \sqrt{5+4x-x^2} dx
  \item \int \frac{t^2}{\sqrt{t^2-6t+13}}dt
\end{mathlist}
\begin{answer}
\begin{mathlist}
  \item \frac{2}{5}
  \item \frac{\pi}{4}+ \frac{\sqrt{3}}{8}-\frac{1}{4}
  \item \ln(\sqrt{x^2+16}+x)-ln(4)+C
  \item \frac{\pi}{16}a^4
  \item \frac{1}{15}\sqrt{t^2+2}(3t^4-8t^2+32)
  \item \frac{81}{4}(\frac{\pi}{8}+ \frac{7}{16} \sqrt{3}-1)
  \item \frac{9}{2}\arcsin \left( \frac{x-2}{3}\right)+ \frac{1}{2}(x-2)\sqrt{5+4x-x^2}+C
  \item \ln\abs{\frac{\sqrt{t^2-6t+13}}{2}+ \frac{t-3}{2}} + C
\end{mathlist}
\end{answer}
\end{exercise}

%------------------------------------------------------------------------------%
\begin{exercise}
Calcule $\int \sin(x) \cos(x)dx$ por quatro métodos.
\begin{alphalist}
  \item Use a substituição $u = \cos(x)$.
  \item Use a substituição $u = \sin(x)$.
  \item Use a identidade $\sin(2x) = 2 \sin(x) \cos(x)$.
  \item Use integração por partes.
\end{alphalist}
\end{exercise}

%------------------------------------------------------------------------------%
\begin{exercise}
Use primeiro uma substituição e depois uma substituição trigonométrica
para calcular a integral
\(\int_0^{\frac{\pi}{2}} \frac{\cos(x)}{\sqrt{1+\sin^2(x)}}dx\).
\end{exercise}

%------------------------------------------------------------------------------%
\begin{exercise}
Use substituição trigonométrica para mostrar que
\[
  \int \frac{dx}{\sqrt{x^2+a^2}}dx= \ln(x+ \sqrt{x^2+a^2})+C.
\]
\end{exercise}

%------------------------------------------------------------------------------%
\begin{exercise}
Calcule o valor médio da função $\frac{\sqrt{x^2-1}}{x}$, $1 \leq x \leq 7$.
\begin{answer}
  $\frac{1}{6}(\sqrt{48}- \arcsec(7))$
\end{answer}
\end{exercise}

%------------------------------------------------------------------------------%
\begin{exercise}
  A parábola $y=\frac{1}{2}x^2$ divide o disco $x^2+y^2 \leq 8$ em duas partes.
  Encontre a área de ambas as partes.
  \begin{answer}
    A área dentro do círculo e acima da parábola é $2\pi +\frac{4}{3}$e a
    área dentro do círculo e abaixo da parábola é $6\pi -\frac{4}{3}$.
  \end{answer}
\end{exercise}
%--------------------------------------------------------------------------------
\begin{exercise}
Verifique se a integral imprópria $\int_{-1}^1 \frac{1}{x^2-2x} dx$ é
convergente e em caso afirmativo calcule seu valor.
\begin{answer}
A integral é divergente.
\end{answer}
\end{exercise}

%------------------------------------------------------------------------------%
%\begin{exercise}
%\begin{mathlist}[2]
%  \item
%  \item
%  \item
%\end{mathlist}
%\begin{answer}
%  \begin{alphalist}[3]
  %    \item
  %    \item
  %    \item
  %  \end{alphalist}
%\end{answer}
%\end{exercise}

\end{exercises}
%------------------------------------------------------------------------------%
